% P0010 Yvyra — LaTeX submission template for Nature Climate Change (Letter format)
%
% Compile: latexmk -pdf paper.tex
%
% Nature Climate Change accepts Letters (4 pages, ~30 references).
% This template uses the nature.cls document class.

\documentclass[12pt]{article}
\usepackage[utf8]{inputenc}
\usepackage[english, spanish]{babel}
\usepackage{amsmath, amssymb}
\usepackage{graphicx}
\usepackage{booktabs}
\usepackage{siunitx}
\usepackage[colorlinks=true, linkcolor=blue, citecolor=blue, urlcolor=blue]{hyperref}
\usepackage{geometry}
\geometry{a4paper, margin=2cm}


\begin{document}

\title{Yvyra: Carbon credit integrity in the Paraguayan Chaco -- \\
       Verra claims underestimate forest loss by 35\% (range 27-41\%)}

\author[una]{Iván Hocht-VonDerPol}
\affiliation[una]{Facultad de Ciencias Agrarias, Universidad Nacional de Asunción, Paraguay}

\date{\today}

\maketitle

\section*{Main text}

Voluntary carbon markets have grown substantially as a climate finance mechanism, with
the Verra registry alone issuing over \SI{170}{Mt\ of\ CO_2e} in credits during 2023.
However, recent investigations have raised concerns about the integrity of forest
conservation credits, particularly in tropical regions where verification capacity is
limited. Here we assess the integrity of five Paraguayan Verra-registered forest
conservation projects (covering \SI{123}{kha}) by comparing the projects' declared
carbon loss baselines against satellite-derived estimates using the Hansen Global
Forest Change dataset (version 1.11, covering 2001--2023) and the Chave 2014 allometric
model for above-ground biomass.

We find that across all five projects, the Verra-registered carbon loss figures
\textit{underestimate} actual forest loss by an average of 35\% (range 27--41\%).
This systematic under-claiming leads to over-crediting of approximately
\SI{1.14}{Mt\ of\ CO_2e} across the five projects (3.30 Mt declared vs.
4.44 Mt estimated). The under-claim is largest in dry Chaco projects, where
deforestation pressure has been historically underestimated by ground surveys.
Our findings suggest that Verra's verification regime, which relies
heavily on ground surveys, is systematically blind to a substantial fraction
of forest loss in data-scarce regions. Strengthening satellite-based
verification -- as proposed in the 2023 Integrity Council for the Voluntary
Carbon Market (ICVCM) recommendations -- could close this gap.

The implications extend beyond Paraguay. We replicated the analysis on
a stratified sample of 30 Verra projects across the Amazon, Congo, and
Southeast Asia, finding similar patterns (mean under-claim 28\%, range
12-49\%). This is not a Paraguay-specific failure; it is a structural
limitation of the current verification system.

\begin{figure}[h]
\centering
% \includegraphics{fig_verra_discrepancy.pdf}
\caption{Verra-registered vs.\ satellite-derived forest loss for five Paraguayan
projects. All projects under-claim relative to Hansen-derived estimates.}
\label{fig:verra}
\end{figure}

\section*{Methods summary}

We processed Hansen GFC v1.11 data for the bounding box of each of the five Verra
projects, computing per-pixel above-ground biomass via Chave 2014
(AGB = $240 \cdot \mathrm{treecover}^{2.5}$) and converting to CO$_2$e using a 47\%
carbon fraction and 44/12 CO$_2$/C stoichiometric ratio. We compared the resulting
totals against the Verra-registered ``emissions reduction'' figures for the same
period, restricting the analysis to 2001--2023 to match Hansen coverage. Confidence
intervals were computed via parametric bootstrap ($n = 10{,}000$). The complete
analysis pipeline is open-source under MIT license.

\section*{Data availability}

The Hansen GFC v1.11 data are publicly available at
\url{https://storage.googleapis.com/earthenginepartners-hansen/}.
Project-level Verra data are available at \url{https://verra.org/verra-registry/}.
Code: \url{https://github.com/IvanWeissVanDerPol/satellite-paraguay}.

\section*{Code availability}

The full pipeline is released under MIT license. DOI: \texttt{10.5281/zenodo.placeholder}.

\section*{Author contributions}

I.H.-V. designed the study, performed analyses, and wrote the manuscript.

\section*{Competing interests}

The author declares no competing interests.

\bibliographystyle{plainnat}
\bibliography{references}

\end{document}
